% --- Archon physics formalization source begin ---
% archon:physics
% archon:covers IPhO2026Problems/problem_ipho_2026_t2_b2.lean
% archon:source-report reports/ipho_2026/problem_ipho_2026_t2_b2.source.json
% archon:problem-id ipho_2026_t2
% archon:part-id T2-B2
% archon:previous-part ipho_2026_t2_b1
% archon:previous-part-policy derive_inline_from_problem_only_material

\chapter{Ipho Physics problem ipho\_2026\_t2\_b2}
\label{ch:IPhO2026Problems_problem_ipho_2026_t2_b2}

\paragraph{Problem source.}
\#\# Physical scenario

A half-cylindrical mirror of radius R illuminates a fully absorbing cylindrical
container of radius a.  Their axes are parallel, and the container center lies
R/2 from the mirror center on the symmetry plane.  Uniform parallel sunlight
arrives along the optical axis.  Any ray absorbed by the container reflects at
most once.  Let theta\_max be the largest incidence angle on the mirror among
rays that strike the container, and let P\_0 be the power the cylinder would
receive without the mirror.  See Figure 2f.

\#\# Current subquestion T2-B2

Express the power ratio P/P\_0 in terms of theta\_max.

\paragraph{Shared problem context.}
A half-cylindrical mirror of radius R illuminates a fully absorbing cylindrical
container of radius a.  Their axes are parallel, and the container center lies
R/2 from the mirror center on the symmetry plane.  Uniform parallel sunlight
arrives along the optical axis.  Any ray absorbed by the container reflects at
most once.  Let theta\_max be the largest incidence angle on the mirror among
rays that strike the container, and let P\_0 be the power the cylinder would
receive without the mirror.  See Figure 2f.

\paragraph{Current subquestion.}
Express the power ratio P/P\_0 in terms of theta\_max.

\paragraph{Figure/image path.}
ipho\_2026\_source/image/T2\_page-3.png

\paragraph{Reusable previous-part conclusions.}
\begin{itemize}
\item Source T2-B1. Question: Given a = alpha*sin(theta\_max) + beta*sin(2*theta\_max), determine alpha and beta in terms of R. Policy: derive\_inline\_from\_problem\_only\_material
\end{itemize}

\paragraph{Formalization target.}
create a compiling Lean file with sorry bodies at `IPhO2026Problems/problem\_ipho\_2026\_t2\_b2.lean`.
The Lean declarations must preserve the physical quantities, dimensions or dimensional roles, figure labels, governing-law hypotheses, and final relation expressed by this problem.
Use Mathlib/Physlib names found through LeanExplore where available. If a domain API is missing, introduce faithful local abstractions rather than scalar placeholder aliases.
This is an answer-blind solve-phase task. Derive a candidate result only from the problem-side evidence above; do not assume a target value, select a tolerance around a desired value, or encode a desired result into a candidate domain.
For numerical questions, define the raw end-to-end quantity and apply an explicit source-derived rounding rule. For identification questions, characterize or prove uniqueness before recording the derived candidate.

\begin{theorem}[Ipho Physics formalization target]
\label{thm:physics:ipho_2026_t2_b2:target}
\lean{IPhO2026.T2B2.power_ratio}
\leanok
The assigned autoformalize agent should translate this IPhO physics problem into Lean declarations in the covered file, with theorem and lemma proof bodies written as `by sorry`.
\end{theorem}
\begin{proof}
\leanok
This is an autoformalization task, not a proof task. Produce faithful statements that can later be proved without weakening the source contract.

\emph{Prover-stage completion.} All 32 proof obligations in
\texttt{IPhO2026Problems/problem\_ipho\_2026\_t2\_b2.lean} are now closed with
kernel-checked proofs (no \texttt{sorry}; only the standard axioms
\texttt{propext}, \texttt{Classical.choice}, \texttt{Quot.sound}).  The shared
T2-B geometry is re-established exactly as in
\texttt{problem\_ipho\_2026\_t2\_b1.lean}: the reflected direction is
$d' = (-\sin 2\theta, -\cos 2\theta)$ (\texttt{reflectedDir\_eq}); the
perpendicular distance from $C = (0, R/2)$ to the reflected line is
$R\sin\theta\,(1-\cos\theta)$ (\texttt{distToLine\_container}), strictly
increasing and continuous on $[0, \pi/2]$ (\texttt{dist\_strictMonoOn},
\texttt{dist\_continuous}); a ray strikes the container iff this distance is at
most $a$ (\texttt{strikes\_iff\_dist\_le}); the extremal ray grazes the
container, $\mathrm{dist}(\theta_{\max}) = a$ (\texttt{dist\_at\_max\_eq}),
giving the factored T2-B1 relation
$a = R\sin\theta_{\max}\,(1-\cos\theta_{\max})$
(\texttt{container\_radius\_eq\_factored}).  For B2 the flux bookkeeping is
added: a reflected ray strikes iff $\theta \le \theta_{\max}$
(\texttt{strikes\_iff\_le\_max}), and $\sin$ covers $[0, \sin\theta_{\max}]$
(\texttt{sin\_surj\_upto}), so the absorbed impact-parameter set is exactly
$[-R\sin\theta_{\max}, R\sin\theta_{\max}]$
(\texttt{absorbedImpactSet\_eq\_Icc}); its Lebesgue width is
$2R\sin\theta_{\max}$ (\texttt{volume\_absorbedImpactSet}, via
\texttt{Real.volume\_Icc}).  Hence per unit length
$P = 2R\sin\theta_{\max}\,I$ (\texttt{powerWithMirror\_eq}) while
$P_0 = 2aI$, so $P/P_0 = R\sin\theta_{\max}/a$
(\texttt{power\_ratio\_eq\_sin}); eliminating $a$ via the factored relation
(with $\sin\theta_{\max} > 0$ and $1-\cos\theta_{\max} > 0$ from
\texttt{thetaMax\_pos}, \texttt{sin\_pos\_of\_theta\_range},
\texttt{one\_sub\_cos\_pos\_of\_theta\_range}) yields
$P/P_0 = 1/(1-\cos\theta_{\max})$ (\texttt{power\_ratio}).
\end{proof}
% --- Archon physics formalization source end ---
